\documentclass[../../../include/open-logic-section]{subfiles}

\begin{document}

\olfileid[id]{sth}{spine}{foundation}

\olsection{Fondasi}

Kita baru \emph{hampir} selesai---dan belum \emph{sepenuhnya}
selesai---karena tidak ada sesuatu pun dalam $\ZFminus$ yang menjamin bahwa
\emph{setiap} himpunan berada dalam suatu $V_\alpha$, yakni bahwa setiap
himpunan dibentuk pada suatu tahap.

Sekarang, terdapat suatu pengertian (matematis) yang cukup langsung yang
menurutnya kita tidak \emph{peduli} apakah terdapat himpunan di luar
hierarki. (Jika memang ada, kita dapat mengabaikannya begitu saja.) Namun,
kita telah memotivasi \emph{konsep} himpunan kita dengan gagasan bahwa setiap
himpunan dibentuk pada suatu tahap (lihat \stageshier{} dalam
\olref[z][story]{sec}). Jadi, kita hendak meniadakan kemungkinan adanya
himpunan yang berada di luar hierarki. Karena itu, kita harus menambahkan
aksioma baru yang memastikan bahwa setiap himpunan muncul di suatu tempat
dalam hierarki.

Karena $V_\alpha$ merupakan tahap-tahap kita, kita dapat saja
mempertimbangkan untuk menambahkan pernyataan berikut sebagai aksioma:

\begin{defish}
\emph{Regularitas.} $\forall A \exists \alpha\, A \subseteq V_\alpha$
\end{defish}

Pendekatan ini sepenuhnya masuk akal. Namun, karena alasan yang akan
dijelaskan pada bagian berikutnya, kita justru akan mengadopsi aksioma
alternatif:

\begin{axiom}[Fondasi]
$(\forall A \neq \emptyset)(\exists B \in A)A \cap B = \emptyset$.
\end{axiom}

Dengan sedikit usaha, kita dapat menunjukkan (dalam $\ZFminus$) bahwa Fondasi
mengimplikasikan Regularitas:
\begin{defn}
Untuk setiap himpunan $A$, tetapkan:
\begin{align*}
	\text{cl}_0(A) &= A,\\
	\text{cl}_{n+1}(A) &= \bigcup \text{cl}_n(A),\\
	\text{trcl}(A) &= \bigcup_{n < \omega} \text{cl}_{n}(A).
\end{align*}
Kita menyebut $\text{trcl}(A)$ sebagai \emph{penutupan transitif} dari $A$.
\end{defn}
\noindent
Nama ``penutupan transitif'' memang tepat:
\begin{prop}\ollabel{subsetoftrcl}
$A \subseteq \trcl{A}$ dan $\trcl{A}$ merupakan himpunan transitif.
\end{prop}

\begin{proof}
Jelas bahwa $A = \text{cl}_0(A) \subseteq \text{trcl}(A)$. Dan jika
$x \in b \in \trcl{A}$, maka $b \in \text{cl}_n(A)$ untuk suatu $n$,
sehingga $x \in \text{cl}_{n+1}(A) \subseteq \text{trcl}(A)$.
\end{proof}

\begin{lem}\ollabel{lem:TransitiveWellFounded}
Jika $A$ merupakan himpunan transitif, maka terdapat suatu $\alpha$ sehingga
$A \subseteq V_\alpha$.
\end{lem}

\begin{proof}
Dengan mengingat definisi ``$\supstrict(X)$'' dari
\olref[ordinals][opps]{defsupstrict}, definisikan dua himpunan:
\begin{align*}
	D  &= \Setabs{x \in A}{\forall \delta\ x \nsubseteq V_\delta}\\
	\alpha &= \supstrict\Setabs{\delta}{(\exists x \in A)
	(x \subseteq V_\delta \land (\forall \gamma \in \delta)x \nsubseteq V_\gamma)}
\end{align*}
Andaikan $D = \emptyset$. Jadi, jika $x \in A$, maka terdapat suatu
$\delta$ sehingga $x \subseteq V_\delta$ dan, berdasarkan pengurutan baik
ordinal, $(\forall \gamma \in \delta)x \nsubseteq V_\gamma$; karena itu
$\delta \in \alpha$, sehingga $x \in V_\alpha$ berdasarkan
\olref[spine][Valphabasic]{Valphabasicprops}. Maka
$A \subseteq V_{\alpha}$, sebagaimana diperlukan.

Jadi, cukup ditunjukkan bahwa $D = \emptyset$. Untuk reductio, andaikan tidak.
Berdasarkan Fondasi, terdapat suatu $B \in D \subseteq A$ sehingga
$D \cap B = \emptyset$. Jika $x \in B$, maka $x \in A$ karena $A$
transitif; dan karena $x \notin D$, diperoleh bahwa
$\exists \delta\ x \subseteq V_\delta$. Sekarang tetapkan
\[
	\beta = \supstrict\Setabs{\delta}{(\exists x \in B)(x \subseteq V_\delta \land (\forall \gamma < \delta)x \nsubseteq V_\gamma)}.
\]
Seperti sebelumnya, $B \subseteq V_\beta$, bertentangan dengan klaim bahwa
$B \in D$.
\end{proof}

\begin{thm}\ollabel{zfentailsregularity}
Regularitas berlaku.
\end{thm}

\begin{proof}
Tetapkan $A$; sekarang $A \subseteq \trcl{A}$ berdasarkan
\olref{subsetoftrcl}, dan yang disebut terakhir bersifat transitif. Jadi,
terdapat suatu
$\alpha$ sehingga $A \subseteq \trcl{A} \subseteq V_\alpha$ berdasarkan
\olref{lem:TransitiveWellFounded}.
\end{proof}

Hasil-hasil ini menunjukkan bahwa $\ZFminus$ membuktikan implikasi
$\emph{Fondasi}\Rightarrow\emph{Regularitas}$. Dalam
\olref[spine][rank]{zfminusregularityfoundation}, kita akan menunjukkan bahwa
$\ZFminus$ membuktikan
$\emph{Regularitas}\Rightarrow\emph{Fondasi}$. Dengan demikian, Fondasi dan
Regularitas \emph{ekuivalen} (modulo $\ZFminus$). Namun, ini berarti bahwa,
dengan mengandaikan $\ZFminus$, kita dapat membenarkan Fondasi dengan
memperhatikan bahwa aksioma itu ekuivalen dengan Regularitas. Dan kita dapat
langsung membenarkan Regularitas berdasarkan \stageshier{}.

\end{document}
